% !TeX root = ../main.tex
\section{Introduction}
\label{sec:introduction}

Rotational plantation management alternates timber harvesting with rapid regeneration and stand growth. Beyond static maps of plantation extent, reliable records of harvest timing, frequency, and inter-harvest intervals are essential for tracking timber supply across rotations~\citep{McMahon2019EucalyptusRotations} and guiding sustainable plantation management~\citep{Cortes2023PlantationHarvesting}. This temporal dimension is particularly critical for Ecuador's commercial balsa (\textit{Ochroma pyramidale}) plantations. Combining exceptionally low wood density with high specific strength, balsa serves as a primary core material in wind-turbine blades and lightweight composite structures~\citep{Borrega2015Balsa_mechanics}. On productive lowland sites, trees reach merchantable size within roughly three years~\citep{CanadasLopez2019BalsaGrowth}, enabling multiple harvest--regrowth cycles on the same parcel within a single decade. Because single-date inventories cannot capture this turnover and sparse observations readily miss harvests obscured by rapid canopy closure~\citep{Huang2026EucalyptusAge}, sustained, high-frequency satellite time series are required to reconstruct rotational harvesting histories.

Long-term optical time-series algorithms such as LandTrendr and CCDC have widely enabled forest disturbance and plantation age reconstruction where clear-sky imagery is abundant~\citep{Kennedy2010LandTrendr,Zhu2014CCDC,Huang2026EucalyptusAge}. Resolving successive short rotations, however, requires uninterrupted observations spanning both canopy removal and intervening regrowth. Across humid tropical regions, persistent cloud cover creates prolonged data gaps~\citep{FloresAnderson2023TropicalCloudFreeOptical}. In mainland Ecuador, year-round cloudiness and cirrus contamination severely restrict usable optical imagery, often forcing regional studies to rely on multi-year optical composites~\citep{OchoaBrito2023EcuadorCropland,AlvarezMendoza2019CirrusEcuador} that sacrifice the monthly precision needed to date individual harvests. Optical archives alone therefore remain insufficient for reconstructing multi-rotation balsa harvest histories.

Sentinel-1 C-band synthetic aperture radar (SAR) overcomes cloud constraints by providing dense, all-weather observations sensitive to canopy geometry and dielectric structure~\citep{Torres2012S1,Reiche2021CongoAlerts}. In short-rotation plantations, clear-felling reduces co- and cross-polarized backscatter, followed by progressive recovery as young stands regenerate~\citep{Spracklen2021Acacia}. Building on this physical response, dense Sentinel-1 time series have supported tropical disturbance alerts, continuous change monitoring, and monthly harvest mapping~\citep{Reiche2021CongoAlerts,Bullock2022CCDC_S1,Zhao2022MonthlyHarvest,Balling2023Texture}.

Extending SAR disturbance detection from isolated alerts to decadal, multi-rotation harvest histories nevertheless faces three distinct challenges. First, across decadal archives, shifting orbital coverage or mixing ascending tracks with different incidence angles and terrain-induced radiometry can introduce step changes unrelated to stand management, compounded by acquisition-level speckle~\citep{Small2011RTC,Quegan2001Speckle}. Second, the canopy baseline itself is non-stationary: repeated harvesting, rapid regrowth, seasonal phenology, and C-band biomass saturation continuously alter the backscatter background across successive rotations~\citep{Santoro2018SARBiomass,Spracklen2021Acacia,VanDerWoude2024ManagementDisturbance}, rendering a single long-term reference inadequate for detecting later harvests. Third, local backscatter shifts do not map one-to-one onto stand-level harvests: rainfall fluctuations can trigger transient backscatter drops without tree removal~\citep{Doblas2020Rainfall}, woody slash can dampen post-harvest contrasts~\citep{Balling2023Texture}, and operational felling often progresses across a stand over several months, generating staggered pixelwise detections from a single harvest episode.

Cumulative sum (CUSUM) analysis provides an effective statistical foundation for isolating sustained backscatter shifts from short-term radar noise by accumulating sequential departures from a reference mean~\citep{Page1954CUSUM,RuizRamos2020CUSUMSMC,Ygorra2021CuSum}. Recent Sentinel-1 adaptations have addressed site-specific tropical forest variability~\citep{Sabir2024AdaptiveCuSUM}, extended detection to multiple disturbances via recursive time-series segmentation in ReCuSum~\citep{Ygorra2023ReCuSum}, and employed short sliding windows for near-real-time alerting in CuSum NRT~\citep{Ygorra2024CuSumNRT}. Yet recursive segmentation makes later detections contingent on earlier split points, whereas short near-real-time windows prioritize low latency at the expense of elevated false alarms, leaving multi-cycle disturbance and return-time reconstruction across decadal SAR archives an open challenge~\citep{Mercier2026SmallLogging}.

Here, we develop and validate a retrospective framework centred on \textit{windowed CuSum} to reconstruct monthly rotational harvest histories from decadal Sentinel-1 observations (May 2016--June 2026) across mainland Ecuador's balsa plantations. Directly addressing the three challenges above, our framework: (1) organizes burst-level RTC observations into track-consistent 30~m spatial time series (STS) to eliminate cross-orbit radiometric shifts while preserving within-stand spatial structure; (2) pairs overlapping context windows with central decision intervals and window-specific permutation testing to renew local baselines across rotations while suppressing boundary and moisture-related false alarms; and (3) consolidates pixelwise detections into stand-level harvest events and intra-event progression groups via temporal detection aggregation (TDA) and spatial timing partitioning (STP). Evaluated against 1,049 dated reference harvests and 13 benchmark detector configurations, the framework enables the first national-scale characterization of harvesting hotspots, market-driven cutting pulses, and climate-linked recurrence intervals across Ecuadorian balsa plantations.
